\chapter{Sample Outputs}

\section{Purpose}
This appendix provides representative outputs generated by the SARS workflow. The examples are formatted to show the structure of extracted academic data, stored analytical results, retrieval chunks, and advisory responses. They illustrate the kind of information that moves between modules in the implemented system.

\section{Representative Marksheet Extraction Output}
\begin{lstlisting}[style=sarscode,caption={Illustrative Extracted Marksheet Structure}]
{
  "hall_ticket_no": "22311A12P9",
  "student_name": "Kethavath Santhosh",
  "branch": "Information Technology",
  "semester_no": 5,
  "exam_month_year": "May 2025",
  "sgpa": 8.41,
  "cgpa": 8.55,
  "total_credits_this_semester": 21.0,
  "subjects": [
    {
      "sno": 1,
      "subject_code": "IT501",
      "subject_name": "Data Warehousing and Mining",
      "grade_letter": "A",
      "grade_points": 8,
      "credits": 3.0,
      "is_backlog": false,
      "result": "P"
    },
    {
      "sno": 2,
      "subject_code": "IT503",
      "subject_name": "Web Technologies",
      "grade_letter": "A+",
      "grade_points": 9,
      "credits": 3.0,
      "is_backlog": false,
      "result": "P"
    }
  ]
}
\end{lstlisting}

\section{Representative Stored Risk Breakdown}
\begin{lstlisting}[style=sarscode,caption={Illustrative Risk Engine Output}]
{
  "gpa_risk": 13.33,
  "backlog_risk": 0.0,
  "attendance_risk": 8.0,
  "gpa_component": 5.33,
  "backlog_component": 0.0,
  "attendance_component": 2.0,
  "sars_score": 7.33,
  "risk_level": "LOW",
  "confidence": 1.0,
  "semesters_analyzed": 5,
  "total_backlogs": 0,
  "cgpa": 8.55,
  "avg_attendance": 69.0,
  "trend_direction": "improving",
  "trend_bonus": -5.0
}
\end{lstlisting}

\section{Representative High-Risk Output}
\begin{lstlisting}[style=sarscode,caption={Illustrative High-Risk Analytical Case}]
{
  "gpa_risk": 44.0,
  "backlog_risk": 90.0,
  "attendance_risk": 22.67,
  "gpa_component": 17.6,
  "backlog_component": 31.5,
  "attendance_component": 5.67,
  "sars_score": 54.77,
  "risk_level": "HIGH",
  "confidence": 0.75,
  "semesters_analyzed": 2,
  "total_backlogs": 4,
  "cgpa": 4.20,
  "avg_attendance": 58.0,
  "trend_direction": "declining_sharply",
  "trend_bonus": 10.0
}
\end{lstlisting}

\section{Representative RAG Chunk}
\begin{lstlisting}[style=sarscode,caption={Illustrative Retrieved Student Chunk}]
{
  "id": "s12_risk_summary",
  "chunk_type": "risk_summary",
  "student_id": "12",
  "text": "Risk summary:
  - SARS score: 82.5
  - Risk level: HIGH
  - Confidence: 0.75
  - Factor breakdown values:
    - cgpa: 4.9
    - total_backlogs: 4
    - trend_direction: declining_sharply
  - Placement eligibility: NOT ELIGIBLE. Current blockers:
    CGPA below 7.5 cutoff, 4 active backlog(s)."
}
\end{lstlisting}

\section{Representative Advisory Interaction}
\begin{lstlisting}[style=sarscode,caption={Illustrative Advisory Exchange}]
Student query:
"What should I focus on first to become placement eligible?"

Retrieved evidence summary:
- Current CGPA is below 7.5
- Student has active backlogs
- Latest SGPA trend is declining

Assistant response:
"Your immediate priority is to clear all active backlog subjects because
any backlog blocks campus placement eligibility. At the same time, your
CGPA must improve toward at least 7.5. Start with backlog subjects that
are core or high-credit, then build a semester plan that prevents further
decline in SGPA."
\end{lstlisting}

\section{Representative Teacher Intervention Record}
\begin{lstlisting}[style=sarscode,caption={Illustrative Stored Intervention Entry}]
{
  "student_id": 12,
  "teacher_id": 3,
  "intervention_type": "Counseling",
  "notes": "Discussed backlog clearance priority and weekly study plan.",
  "status": "open",
  "created_at": "2026-04-08T18:35:14Z"
}
\end{lstlisting}

\section{Interpretive Value of the Outputs}
The sample outputs show that the project does not only display information in the UI. It maintains a chain of machine-readable intermediate artifacts:
\begin{itemize}[leftmargin=*]
\item extracted semester structures,
\item validated academic records,
\item factor-wise risk breakdowns,
\item retrieved knowledge chunks,
\item intervention records and advisory conversations.
\end{itemize}

These artifacts are valuable because they make the system inspectable. A reviewer can trace how the project moves from document ingestion to advisory response through stable intermediate representations.
